% GENERATED FILE. Edit canonical lessons, then run npm run generate:books.
% canonical-source-hash: fnv1a64:47f3f3f7c9810554
% canonical-lessons: KA-C290-ardha, KA-C290-purti, KA-C290-kasta, KA-C290-sulabha, KA-C290-mukhya

\chapter{Half, Whole, Difficult, Easy, Important}
\label{ch:ka-half-whole-difficult-easy-important}
\hlchaptermodality{290}

\begin{chapteropening}
\textbf{By the end of this chapter:} \emph{I can say half, whole, difficult, easy and important.}

Complete the last lesson of chapter 290: I can say half, whole, difficult, easy and important.
\end{chapteropening}

\section[\texorpdfstring{\kn{ಅರ್ಧ} (ardha)}{ಅರ್ಧ (ardha)}]{\kn{ಅರ್ಧ} (ardha) — half}

\label{lesson:KA-C290-ardha}



\begin{quote}
Before the new one: say the Kannada for a basis, then the Kannada for evidence.
\end{quote}

\subsection*{You'll want to know: \kn{ಅರ್ಧ}}
\textbf{\kn{ಅರ್ಧ}} — \emph{ardha} — \textquotedblleft{}half\textquotedblright{}.

\begin{grammarlens}[title={What you've built}]
One more word for this chapter.
\end{grammarlens}

\subsection*{Your turn}
\begin{itemize}
\raggedright
  \item \emph{Say it:} \emph{ardha}
  \item \emph{Say it:} \emph{ardha}, once more
  \item \emph{Say it:} say \emph{ardha}
  \item \emph{Recall:} say the Kannada for a basis, then the Kannada for evidence, then say \emph{ardha} again
\end{itemize}

\begin{tcolorbox}[breakable,colback=teal!4,colframe=teal!35!black,title={Before you move on}]
What does \kn{ಅರ್ಧ} mean? (\textquotedblleft{}Half\textquotedblright{}.) Say it once more.
\end{tcolorbox}
\section[\texorpdfstring{\kn{ಪೂರ್ತಿ} (pūrti)}{ಪೂರ್ತಿ (pūrti)}]{\kn{ಪೂರ್ತಿ} (pūrti) — whole, complete}

\label{lesson:KA-C290-purti}



\begin{quote}
Before the new one: say the Kannada for evidence, then the Kannada for half.
\end{quote}

\subsection*{You'll want to know: \kn{ಪೂರ್ತಿ}}
\textbf{\kn{ಪೂರ್ತಿ}} — \emph{pūrti} — \textquotedblleft{}whole, complete\textquotedblright{}.

\begin{grammarlens}[title={What you've built}]
One more word for this chapter.
\end{grammarlens}

\subsection*{Your turn}
\begin{itemize}
\raggedright
  \item \emph{Say it:} \emph{pūrti}
  \item \emph{Say it:} \emph{pūrti}, once more
  \item \emph{Say it:} say \emph{pūrti}
  \item \emph{Recall:} say the Kannada for evidence, then the Kannada for half, then say \emph{pūrti} again
\end{itemize}

\begin{tcolorbox}[breakable,colback=teal!4,colframe=teal!35!black,title={Before you move on}]
What does \kn{ಪೂರ್ತಿ} mean? (\textquotedblleft{}Whole, complete\textquotedblright{}.) Say it once more.
\end{tcolorbox}
\section[\texorpdfstring{\kn{ಕಷ್ಟ} (kaṣṭa)}{ಕಷ್ಟ (kaṣṭa)}]{\kn{ಕಷ್ಟ} (kaṣṭa) — difficult}

\label{lesson:KA-C290-kasta}



\begin{quote}
Before the new one: say the Kannada for half, then the Kannada for whole.
\end{quote}

\subsection*{You'll want to know: \kn{ಕಷ್ಟ}}
\textbf{\kn{ಕಷ್ಟ}} — \emph{kaṣṭa} — \textquotedblleft{}difficult\textquotedblright{}.

\begin{grammarlens}[title={What you've built}]
One more word for this chapter.
\end{grammarlens}

\subsection*{Your turn}
\begin{itemize}
\raggedright
  \item \emph{Say it:} \emph{kaṣṭa}
  \item \emph{Say it:} \emph{kaṣṭa}, once more
  \item \emph{Say it:} say \emph{kaṣṭa}
  \item \emph{Recall:} say the Kannada for half, then the Kannada for whole, then say \emph{kaṣṭa} again
\end{itemize}

\begin{tcolorbox}[breakable,colback=teal!4,colframe=teal!35!black,title={Before you move on}]
What does \kn{ಕಷ್ಟ} mean? (\textquotedblleft{}Difficult\textquotedblright{}.) Say it once more.
\end{tcolorbox}
\section[\texorpdfstring{\kn{ಸುಲಭ} (sulabha)}{ಸುಲಭ (sulabha)}]{\kn{ಸುಲಭ} (sulabha) — easy}

\label{lesson:KA-C290-sulabha}



\begin{quote}
Before the new one: say the Kannada for whole, then the Kannada for difficult.
\end{quote}

\subsection*{You'll want to know: \kn{ಸುಲಭ}}
\textbf{\kn{ಸುಲಭ}} — \emph{sulabha} — \textquotedblleft{}easy\textquotedblright{}.

\begin{grammarlens}[title={What you've built}]
One more word for this chapter.
\end{grammarlens}

\subsection*{Your turn}
\begin{itemize}
\raggedright
  \item \emph{Say it:} \emph{sulabha}
  \item \emph{Say it:} \emph{sulabha}, once more
  \item \emph{Say it:} say \emph{sulabha}
  \item \emph{Recall:} say the Kannada for whole, then the Kannada for difficult, then say \emph{sulabha} again
\end{itemize}

\begin{tcolorbox}[breakable,colback=teal!4,colframe=teal!35!black,title={Before you move on}]
What does \kn{ಸುಲಭ} mean? (\textquotedblleft{}Easy\textquotedblright{}.) Say it once more.
\end{tcolorbox}
\section[\texorpdfstring{\kn{ಮುಖ್ಯ} (mukhya)}{ಮುಖ್ಯ (mukhya)}]{\kn{ಮುಖ್ಯ} (mukhya) — important}

\label{lesson:KA-C290-mukhya}



\begin{quote}
Before the new one: say the Kannada for difficult, then the Kannada for easy.
\end{quote}

\subsection*{You'll want to know: \kn{ಮುಖ್ಯ}}
\textbf{\kn{ಮುಖ್ಯ}} — \emph{mukhya} — \textquotedblleft{}important\textquotedblright{}.

\begin{grammarlens}[title={What you've built}]
One more word for this chapter.
\end{grammarlens}

\subsection*{Your turn}
\begin{itemize}
\raggedright
  \item \emph{Say it:} \emph{mukhya}
  \item \emph{Say it:} \emph{mukhya}, once more
  \item \emph{Say it:} say \emph{mukhya}
  \item \emph{Recall:} say the Kannada for difficult, then the Kannada for easy, then say \emph{mukhya} again
\end{itemize}

\begin{tcolorbox}[breakable,colback=teal!4,colframe=teal!35!black,title={Before you move on}]
What does \kn{ಮುಖ್ಯ} mean? (\textquotedblleft{}Important\textquotedblright{}.) Say it once more.
\end{tcolorbox}
